\FloatBarrier
\section{From Orbital Hamiltonians to Spin Exchange}
\label{app:exchange-reduction}
\label{sec:effective-models-numerics}

The general many-electron matrix in
Appendix~\ref{sec:device-orbital-basis} retains all interactions within
a chosen orbital space. This appendix develops the more restrictive
model used to derive the exchange interaction in
Section~\ref{sec:exchange-virtual-tunneling}. We first specify the
conditions for a single-spatial-state Hubbard model. We then eliminate
its high-energy charge states to obtain spin exchange, check the
splitting against exact diagonalization, and calculate the singlet
charge admixture that enters interactions between encoded qubits.

\subsection{One spatial state per dot: assumptions of the Hubbard model}
\label{sec:device-hubbard-reduction}

For two neighboring spin qubits, the relevant charge states include
one electron on each dot and the doubly occupied states reached by
moving an electron between them. The one-electron spatial states are
concentrated near individual dots, but their tails extend into
neighboring wells and permit tunneling. Repulsion changes the energy
cost of this transfer. We need a model that retains these charge
states and the tunneling between them before deriving their effect
on the spins. Starting from Appendix~\ref{sec:device-orbital-basis},
we construct such a model for dots with a small number of electrons.

\paragraph{Restricting the spatial states.}
Choose one orthonormal localized spatial state $w_i$ per dot,
including its local valley composition, and retain its two spin
states. We neglect spin-flip tunneling. The retained spatial state
has energy $\epsilon_i$ and can hold zero, one, or two electrons;
the doubly occupied state is a spin singlet. We use this same spatial
basis for the charge configurations being compared. The model
therefore retains changes in local occupation, including double
occupation during tunneling.

This restriction assumes that other orbital--valley states have
negligible thermal occupation and that their mixing into the retained
states gives corrections below the required accuracy. The condition
must hold in every charge sector used in the calculation, including
the doubly occupied states that enter a virtual tunneling process.
A large single-electron excitation gap alone does not establish this:
electron repulsion can change the ordering and composition of the
two-electron states. If excited orbitals or valleys contribute
appreciably, the local spatial basis must be enlarged.

Putting two opposite-spin electrons on the same dot costs an
additional repulsion energy $U_i$. Occupying different dots $i,j$
gives a direct repulsion $V_{ij}$. Write $|w_i(\mathbf r)|^2$
for the spatial probability density, summed over retained valley
components when needed. From the pair matrix elements,
\begin{equation}
 \begin{aligned}
 U_i&=\mathcal V_{iiii}
   =\int d^3r\,d^3r'\,
       |w_i(\mathbf r)|^2W(\mathbf r,\mathbf r')
                         |w_i(\mathbf r')|^2,\\
 V_{ij}&=\mathcal V_{ijij}
   =\int d^3r\,d^3r'\,
       |w_i(\mathbf r)|^2W(\mathbf r,\mathbf r')
                         |w_j(\mathbf r')|^2.
 \end{aligned}
 \label{eq:device-density-interactions}
\end{equation}
These integrals calculate the repulsion energies from the
orbital charge distributions and their screened interaction.

\paragraph{Restricting the interaction terms.}
Within this spatial basis, a further approximation retains only
$U_i$ and $V_{ij}$ from the pair interaction. This requires the
remaining pair matrix elements to give negligible corrections to
the energies or dynamics being calculated. The spatial-state
truncation alone does not imply this condition.
Keeping these density interactions together with the one-electron
tunneling from Eq.~\eqref{eq:device-one-body-matrix} gives the
\emph{extended Hubbard Hamiltonian}
\begin{equation}
 \begin{aligned}
 H_{\mathrm{Hub}}
   &=\sum_i\epsilon_i n_i
     -\sum_{i<j,\sigma}
        \left(t_{ij}c_{i\sigma}^\dagger c_{j\sigma}
             +t_{ij}^*c_{j\sigma}^\dagger c_{i\sigma}\right)\\
   &\quad+\sum_i U_i n_{i\uparrow}n_{i\downarrow}
          +\sum_{i<j}V_{ij}n_i n_j,\\
 n_{i\sigma}&=c_{i\sigma}^\dagger c_{i\sigma},
 \qquad n_i=n_{i\uparrow}+n_{i\downarrow}.
 \end{aligned}
 \label{eq:device-extended-hubbard}
\end{equation}
The operator $n_i$ counts electrons on dot $i$.
Each displayed term now has a specific origin: localized-state
energy, coherent transfer, same-dot repulsion, and interdot
repulsion. Local Zeeman terms can be added when resolving spins.

On a charge configuration with $N_i=0,1,2$, the operators $n_i$
and $n_{i\uparrow}n_{i\downarrow}$ have eigenvalues $N_i$ and
$N_i(N_i-1)/2$, respectively. In the occupation-number basis of
this Hubbard model, the diagonal matrix element is therefore
\begin{equation}
 E_{\mathrm{ch}}(\{N_i\})
   =\sum_i\epsilon_i N_i
    +\sum_i U_i\frac{N_i(N_i-1)}2
    +\sum_{i<j}V_{ij}N_iN_j.
 \label{eq:dot-array-charge-energy}
\end{equation}
Tunneling supplies off-diagonal matrix elements between these basis
states, so the many-electron eigenenergies require diagonalizing
$H_{\mathrm{Hub}}$. The expression above relies on the spatial-state
and interaction restrictions already specified; spin-dependent
energies are omitted here.

This model omits other allowed pair matrix elements. For example,
the Coulomb interaction connects the spatial product states
$w_i(\mathbf r_1)w_j(\mathbf r_2)$ and
$w_j(\mathbf r_1)w_i(\mathbf r_2)$, with matrix element
$K_{ij}=\mathcal V_{ijji}$. In the subspace with one electron
in each orbital and no tunneling, their symmetric spatial
combination has energy $V_{ij}+K_{ij}$, and their antisymmetric
combination has energy $V_{ij}-K_{ij}$.
Fermionic antisymmetry pairs these with a spin singlet and triplet,
respectively. This \emph{direct Coulomb exchange} therefore contributes
$E_T-E_S=-2K_{ij}$ through the exchange symmetry of the spatial
state at zero tunneling. When calculating a small spin-exchange
splitting, the size of $2K_{ij}$ must be judged against that splitting
and the accuracy sought; its being small compared with $U_i$ is
insufficient.

The full interaction can also transfer two electrons together
between orbitals, or transfer one with an amplitude depending on
the other electron's orbital. These are called pair hopping and
occupation-dependent hopping~\cite{Wang2011}. They can be restored
from $\mathcal V_{ijkl}$ if their effects matter.

The Hubbard model obtained under these conditions still contains
empty and doubly occupied dots, and allows tunneling to mix charge
configurations. A Hamiltonian acting only on spins requires the
further charge-state elimination in
Appendix~\ref{sec:device-low-energy-reduction}: there we restrict to
one electron per dot and require the positive double-occupation
costs to be much larger than the tunnel coupling.

\subsection{Keeping spin states and eliminating costly excitations}
\label{sec:device-low-energy-reduction}

Starting from the Hubbard Hamiltonian in
Eq.~\eqref{eq:device-extended-hubbard}, we seek a description in
which the electron number and lowest orbital--valley state on each
dot are fixed, while spins remain free to evolve.
Virtual tunneling through excited charge states shifts the
energies within this low-energy set. We retain these shifts
when deriving its effective Hamiltonian.

To derive their effect, let $P$ project onto the states we want
to keep, and $Q=I-P$ onto the remaining states, with $I$ the identity.
Write an eigenstate as
$\ket{\Psi}=\ket p+\ket q$, where
$\ket p=P\ket\Psi$ and $\ket q=Q\ket\Psi$.
Its two coupled equations are
\begin{equation}
 \begin{aligned}
 PHP\ket p+PHQ\ket q&=E\ket p,\\
 QHP\ket p+QHQ\ket q&=E\ket q.
 \end{aligned}
 \label{eq:device-block-eigenproblem}
\end{equation}
Solving the second for $\ket q$ and substituting into the first
gives
\begin{equation}
 \begin{aligned}
 \ket q&=(E-QHQ)^{-1}QHP\ket p,\\
 H_{\mathrm{eff}}(E)&=PHP+
                PHQ(E-QHQ)^{-1}QHP.
 \end{aligned}
 \label{eq:device-effective-resolvent}
\end{equation}
The inverse is taken only in the excluded space, away from its
eigenvalues. Reading the second term from right to left,
$QHP$ takes a retained state out, the energy denominator weights
the intermediate amplitude, and $PHQ$ brings it back.
This is the origin of a virtual correction.

If the coupling between the two subspaces is small compared
with their energy separation, this expression can be expanded.
An equivalent construction, the Schrieffer--Wolff transformation,
produces an energy-independent low-energy Hamiltonian
order by order~\cite{Bravyi2011}. With one electron per dot
and approximately spin-rotation-invariant exchange, its leading
form is
\begin{equation}
 H_{\mathrm{spin}}=E_{\mathrm{off}}
       +\sum_i\mathbf b_i\cdot\mathbf s_i
       +\sum_{i<j}J_{ij}\mathbf s_i\cdot\mathbf s_j
       +\delta H_{\mathrm{spin}},
 \qquad \mathbf s_i=\boldsymbol\sigma_i/2.
 \label{eq:device-spin-hamiltonian}
\end{equation}
$E_{\mathrm{off}}$ is a common energy shift,
$\mathbf b_i$ is the local Zeeman energy vector, and
$\boldsymbol\sigma_i$ is the vector of Pauli operators for spin $i$.
The remainder $\delta H_{\mathrm{spin}}$ contains additional terms,
such as anisotropic exchange and multispin interactions.

For two dots, a hop from the $(1,1)$ configuration produces
$(2,0)$ or $(0,2)$. These states are intermediate because
a second hop returns to $(1,1)$.
With detuning $\varepsilon=\epsilon_1-\epsilon_2$, their costs
and leading tunneling-induced exchange are
\begin{equation}
 \begin{aligned}
 \Delta_{20}&=U_1-V_{12}+\varepsilon,
 &\Delta_{02}&=U_2-V_{12}-\varepsilon,\\
 J_{\mathrm{se}}&=2|t_{12}|^2
       \left(\frac1{\Delta_{20}}+\frac1{\Delta_{02}}\right).
 \end{aligned}
 \label{eq:device-general-superexchange}
\end{equation}
Both costs must be positive and much larger than $|t_{12}|$.
A small direct-exchange correction contributes $-2K_{12}$ at
leading order. If a charge, orbital, or valley excitation becomes
comparable to its coupling to the retained states, it should be
kept explicitly in the retained space.

\paragraph{Fermionic matrix elements and charge elimination.}
The exchange result used in Eq.~\eqref{eq:superexchange} follows
by retaining the singlet's two virtual charge channels. Specialize
Eq.~\eqref{eq:device-extended-hubbard} to two dots with
$U_1=U_2=U$, $V_{12}=V$, and real $t_{12}=t_c$.
Write $H_t$ for its tunneling term.
The charge costs $\Delta_{20},\Delta_{02}$ are given in
Eq.~\eqref{eq:virtual-charge-gaps}.
We evaluate the tunneling matrix elements in the $(1,1)$
singlet--triplet basis
\begin{equation}
  \begin{aligned}
    \ket{S_{11}}&=\frac{\ket{\uparrow\downarrow}
                          -\ket{\downarrow\uparrow}}{\sqrt{2}},&
    \ket{T_{0,11}}&=\frac{\ket{\uparrow\downarrow}
                            +\ket{\downarrow\uparrow}}{\sqrt{2}},\\
    \ket{T_{+,11}}&=\ket{\uparrow\uparrow},&
    \ket{T_{-,11}}&=\ket{\downarrow\downarrow}.
  \end{aligned}
  \label{eq:hubbard-spin-states}
\end{equation}
The subscript $11$ labels charge occupation; $\ket{T_{11}}$ below
denotes any triplet. Fix the doubly occupied singlet phases by
$\ket{D_{20}}=c_{1\uparrow}^\dagger c_{1\downarrow}^\dagger
\ket{\mathrm{vac}}$ and
$\ket{D_{02}}=c_{2\uparrow}^\dagger c_{2\downarrow}^\dagger
\ket{\mathrm{vac}}$, where $\ket{\mathrm{vac}}$ is the empty-dot state.
For product states, create the dot-1 electron before the dot-2
electron, so
$\ket{\uparrow\downarrow}=c_{1\uparrow}^\dagger
c_{2\downarrow}^\dagger\ket{\mathrm{vac}}$. Applying $H_t$ gives
\begin{equation}
  \begin{aligned}
    H_t\ket{\uparrow\downarrow}
       &=-t_c(\ket{D_{20}}+\ket{D_{02}}),\\
    H_t\ket{\downarrow\uparrow}
       &=+t_c(\ket{D_{20}}+\ket{D_{02}}).
  \end{aligned}
  \label{eq:tunneling-product-states}
\end{equation}
The two amplitudes add for $\ket{S_{11}}$ and cancel for
$\ket{T_{0,11}}$. Pauli exclusion blocks the polarized triplets.
Thus the singlet couples to both doubly occupied states, while
every triplet is uncoupled:
\begin{equation}
  |\bra{D_{20}}H_t\ket{S_{11}}|
   =|\bra{D_{02}}H_t\ket{S_{11}}|=\sqrt{2}|t_c|,\qquad
  \bra{D_{20}}H_t\ket{T_{11}}=
  \bra{D_{02}}H_t\ket{T_{11}}=0.
  \label{eq:tunnel-singlet-matrix-elements}
\end{equation}

\paragraph{Eliminating charge excitations.}
We now calculate the singlet energy shift caused by its two
tunneling channels. In the basis
$(\ket{S_{11}},\ket{D_{20}},\ket{D_{02}})$, subtracting
the reference energy $E_{11}=\epsilon_1+\epsilon_2+V$ gives the matrix
\begin{equation}
  M_S=\begin{pmatrix}
       0&-\sqrt{2}t_c&-\sqrt{2}t_c\\
       -\sqrt{2}t_c&\Delta_{20}&0\\
       -\sqrt{2}t_c&0&\Delta_{02}
      \end{pmatrix}.
  \label{eq:singlet-charge-matrix}
\end{equation}
Let
$\ket{\widetilde S}=a_S\ket{S_{11}}+a_{20}\ket{D_{20}}
+a_{02}\ket{D_{02}}$ be its normalized low-energy eigenstate,
with energy shift $\delta E_S$.
To eliminate the excited charge states, solve for their amplitudes
$a_{20},a_{02}$ in terms of $a_S$. The eigenvalue equation reads
\begin{equation}
  \begin{aligned}
    \delta E_S a_S&=-\sqrt{2}t_c(a_{20}+a_{02}),\\
    (\Delta_{20}-\delta E_S)a_{20}&=\sqrt{2}t_c a_S,\\
    (\Delta_{02}-\delta E_S)a_{02}&=\sqrt{2}t_c a_S.
  \end{aligned}
  \label{eq:charge-amplitude-equations}
\end{equation}
For $|\delta E_S|\ll\Delta_{20},\Delta_{02}$, the last two equations give
\begin{equation}
  a_{20}\simeq\frac{\sqrt{2}t_c}{\Delta_{20}}a_S,\qquad
  a_{02}\simeq\frac{\sqrt{2}t_c}{\Delta_{02}}a_S.
  \label{eq:virtual-charge-amplitudes}
\end{equation}
Substitution into the first line of
Eq.~\eqref{eq:charge-amplitude-equations} gives the singlet shift;
the triplet has no corresponding channel:
\begin{equation}
  \delta E_S=-\frac{2t_c^2}{\Delta_{20}}
              -\frac{2t_c^2}{\Delta_{02}},\qquad
  \delta E_T=0.
  \label{eq:singlet-triplet-virtual-shifts}
\end{equation}
Their difference is the exchange energy $J=E_T-E_S$ in
Eq.~\eqref{eq:superexchange}. The correction within $(1,1)$ is
$-JP_S$, where $P_S=\ket{S_{11}}\bra{S_{11}}$.
Since $\mathbf s_1\cdot\mathbf s_2$ has singlet and triplet
eigenvalues $-3/4$ and $1/4$,
$P_S=I/4-\mathbf s_1\cdot\mathbf s_2$.
Dropping the common energy shift $-J/4$ gives
$H_{\mathrm{ex}}=J\mathbf s_1\cdot\mathbf s_2$.
Squaring the small charge amplitudes, with $|a_S|^2\simeq1$,
gives the total double-occupation weight in the low-energy singlet:
\begin{equation}
  P_D=|a_{20}|^2+|a_{02}|^2
       \simeq 2t_c^2\left(\frac{1}{\Delta_{20}^2}
                          +\frac{1}{\Delta_{02}^2}\right)\ll1.
  \label{eq:double-occupation-weight}
\end{equation}

\paragraph{Comparison with exact diagonalization.}
We can check this low-energy expansion against exact
diagonalization of the two-site Hubbard Hamiltonian.
Take $U_1=U_2=U$, $V_{12}=V$, zero detuning,
and real hopping $t_{12}=t$, with Zeeman and direct-exchange
terms omitted. Set $D=U-V>0$. With two electrons, tunneling
couples the $(1,1)$ singlet $\ket{S_{11}}$ to the symmetric
combination of the doubly occupied states,
$\ket{D_+}=(\ket{\uparrow\downarrow,0}
+\ket{0,\uparrow\downarrow})/\sqrt2$.
Relative to the triplet energy, the matrix in the basis
$(\ket{S_{11}},\ket{D_+})$ and its exact singlet--triplet
splitting are
\begin{equation}
 M_S=\begin{pmatrix}0&-2t\\-2t&D\end{pmatrix},
 \qquad
 J_{\mathrm{exact}}\equiv E_T-E_S
     =\frac{\sqrt{D^2+16t^2}-D}{2}.
 \label{eq:device-exact-hubbard-check}
\end{equation}
Diagonalization retains double occupation explicitly. Expanding
for $|t|\ll D$ gives $J_{\mathrm{exact}}\simeq4t^2/D$, agreeing
with Eq.~\eqref{eq:device-general-superexchange} when
$\Delta_{20}=\Delta_{02}=D$.

Virtual processes can also involve more than two dots.
On a four-dot loop, four successive hops can restore one
electron per dot while cyclically permuting the spins.
For comparable hopping amplitudes $t$ and charge gaps $\Delta$,
the resulting cyclic exchange is of order $t^4/\Delta^3$.
The effective multispin term thus arises from successive virtual
transfers generated by one-electron tunneling and pairwise
repulsion.

Finally, a small model can be constructed either by evaluating
matrix elements in a chosen basis or by adjusting its parameters
to reproduce selected energies from a larger calculation.
The second procedure can incorporate excluded-state effects
into effective parameters. Matching theoretical energies in
this way is a theoretical model reduction. Its accuracy must be tested in the states and
voltage range where the small model will be used.
